\documentclass[letterpaper,9pt]{extarticle}

\usepackage{latexsym}
\usepackage[empty]{fullpage}
\usepackage{titlesec}
\usepackage{marvosym}
\usepackage[usenames,dvipsnames]{color}
\usepackage{verbatim}
\usepackage{enumitem}
\usepackage[pdftex]{hyperref}
\usepackage{fancyhdr}
\usepackage{fontawesome5}

\pagestyle{fancy}
\fancyhf{} % clear all header and footer fields
\fancyfoot{}
\renewcommand{\headrulewidth}{0pt}
\renewcommand{\footrulewidth}{0pt}

\hypersetup{hidelinks}

% Adjust margins
\addtolength{\oddsidemargin}{-0.50in}
\addtolength{\evensidemargin}{-0.40in}
\addtolength{\textwidth}{1in}
\addtolength{\topmargin}{-.7in}
\addtolength{\textheight}{1.1in}


\urlstyle{same}

\raggedbottom
\raggedright
\setlength{\tabcolsep}{0in}

% Sections formatting
\titleformat{\section}{
  \vspace{-7pt}\scshape\raggedright\large
}{}{0em}{}[\color{black}\titlerule \vspace{-5pt}]

%-------------------------
% Custom commands


% \resumeItem — bold header followed by a nested bulleted list (used inside Experience clusters)
\newcommand{\resumeItem}[2]{
  \item\small{
    \textbf{#1}{:} 
  \begin{itemize}[label=\textbullet, leftmargin=8pt, itemsep=2pt]
      #2
    \end{itemize}
    }
}

% \resumeSubItem — single-line "Title: description" entry (used in Projects, Skills, Accolades)
\newcommand{\resumeSubItem}[2]{
  \item\small{
    \textbf{#1}{: #2}
  }\vspace{-2pt}
}

% \resumeSubheading — job/education header: Company/Location on top row, Role/Dates on bottom row
\newcommand{\resumeSubheading}[4]{
  \item
    \begin{tabular*}{0.97\textwidth}{l@{\extracolsep{\fill}}r}
      \textbf{#1} & {\small #2} \\
      {\small#3} & {\small #4} \\
    \end{tabular*}
}

\renewcommand{\labelitemii}{$\circ$}

\newcommand{\resumeSubHeadingListStart}{\begin{itemize}[leftmargin=*]}
\newcommand{\resumeSubHeadingListEnd}{\end{itemize}}
\newcommand{\resumeItemListStart}{\begin{itemize}[leftmargin=6pt, itemindent=0pt, labelsep=4pt]}
\newcommand{\resumeItemListEnd}{\end{itemize}}

%-------------------------------------------
%%%%%%  CV STARTS HERE  %%%%%%%%%%%%%%%%%%%%%%%%%%%%


\begin{document}

%----------HEADING-----------------

\begin{tabular*}{\textwidth}{l@{\extracolsep{\fill}}r}

    \textbf{{\large Rituraj Kulshresth}} & Email:\href{mailto:kulshresth.1@alumni.iitj.ac.in}{ kulshresth.1@alumni.iitj.ac.in}\\
    
     \href{https://github.com/RiturajKulshresth}{GitHub} \textbar{} \href{https://riturajkulshresth.vercel.app}{Portfolio} & Location: Hyderabad, India \\
    
     \href{https://www.linkedin.com/in/rituraj-kulshresth/}{LinkedIn}& Mobile: +91 7004742004 \\
  
\end{tabular*}

%-----------SUMMARY-----------------
\section{Summary}
\small{Software engineer specializing in enterprise AI platforms. Designed and shipped three production agent and RAG platforms at Warner Bros. Discovery used by 11+ internal teams. Architecture lead on organization-wide documents for agent authentication, evaluation, and retrieval.}

%-----------EDUCATION-----------------

\section{Education}
  \resumeSubHeadingListStart
    \resumeSubheading
      {Indian Institute of Technology, Jodhpur}{Jodhpur, Rajasthan}
      {Bachelor of Technology in Computer Science and Engineering;}{July 2018 -- June 2022}
  \resumeSubHeadingListEnd

%-----------EXPERIENCE-----------------
\section{Experience}
  \resumeSubHeadingListStart

    \resumeSubheading
          {Warner Bros. Discovery}{Hyderabad, India}
          {Software Development Engineer 2 (AI Platforms)}{Nov 2024 - Present}
          \resumeItemListStart
            \resumeItem{ACME (Agent Core for Managed Execution): Enterprise Agent Platform}
            {
                \item Designed ACME's credentials and authentication architecture: a \textbf{four-resource model} (Team, App, Agent, Sub-agent) with separate Edit and Invocation ACLs and dual auth modes (\texttt{act\_as\_app} via API keys and AWS Secrets Manager; \texttt{act\_as\_user} via Okta OIDC), plus a \textbf{sub-agent safety property} propagating only the caller's identity through every tool call, making cross-team agent reuse safe by construction.

                \item Built the \textbf{ACME Playground}, a sandboxed surface to prompt-test agents, inspect tool calls and streaming events, and validate sub-agent orchestration before publishing for org-wide reuse.

                \item Refactored agent configuration off local YAML into a versioned \textbf{PostgreSQL (JSONB) registry} with a draft/published/archived lifecycle and API-driven Console Builder and Catalog, enabling \textbf{multi-tenant onboarding} without code changes.

                \item Shipped \textbf{multi-model serving} via AWS Bedrock inference profiles alongside OpenAI and Anthropic, standardizing model routing and fallback for production traffic.

                \item Drove ACME from \textbf{CRITICAL}-severity SAST findings to \textbf{zero} in a single sprint, hardening release posture ahead of org rollout.

                \item Provisioned ACME infrastructure on the org IaC platform (Terraform/Terragrunt, Helm) and bootstrapped the \textbf{BOLT dual-language monorepo} (FastAPI + Next.js) with CI/CD across DEV, INT, STG, and PROD.
            }
            \vspace{2pt}

            \resumeItem{Knowledge Hub: Enterprise RAG-as-a-Service}
            {
                \item Built Knowledge Hub \textbf{end-to-end}: a \textbf{self-serve retrieval service} on \textbf{AWS OpenSearch} where teams create enterprise document indices and share document and index catalogs org-wide, so any AI app can reuse existing indices without re-ingesting; consumed by multiple internal apps and backed by an org-wide architecture document I authored.

                \item Implemented ingestion from Confluence, Jira, GitHub, S3, Airflow logs, and arbitrary text, with bulk upload, team-specific chunking, and a \textbf{semantic chunking} mode that splits by meaning instead of token windows.

                \item Added \textbf{multimodal embedding support} (Bedrock Titan and Anthropic Claude for text/code, a Bedrock multimodal model for images), extending access to teams whose documents are diagrams, screenshots, and scanned PDFs.
            }
            \vspace{2pt}

            \resumeItem{Daisy: Data Analytics AI Platform}
            {
                \item Developed Daisy, a \textbf{multi-agent AI platform} used by \textbf{11+} engineering and business teams to build domain-specific chatbots with fine-grained control over tools, data sources (SQL, Text), and models, using Vercel AI SDK, FastAPI, and LangChain.

                \item Cut Daisy infra spend from \textbf{$\sim$\$700 to $\sim$\$200 per environment} (DEV/INT/STG) and \textbf{$\sim$\$1000 to $\sim$\$800 in PROD} by right-sizing compute/storage and decommissioning unused services.

                \item Built an evaluation framework using DeepEval and custom metrics for prompt quality and LLM correctness, improving evaluation accuracy by \textbf{27\%}.

                \item Cut chat-session storage by \textbf{$\sim$66\%} by optimizing Vercel AI SDK logging, removing redundancy and improving OpenSearch indexing and query performance.

                \item Led migration of the entire Daisy platform to the org's new in-house IaC accounts during an org split: moved codebases and Lambdas off the legacy Discovery account and reprovisioned infrastructure with Terraform/Terragrunt, maintaining \textbf{zero downtime}.

                \item Authored \textbf{org-wide Architecture Documents} (Evaluation Platform, Daisy Vector DB API) and ran code reviews for a team of \textbf{7 engineers}.
            }
        \resumeItemListEnd
        \vspace{3pt}

    \resumeSubheading
        {Deloitte}{Hyderabad, India}
        {Software Engineer (Analyst)}{July 2022 -- Nov 2024}
        \resumeItemListStart
            \resumeItem{Outlook add-ins and Document Archival Apps}
            {
                \item Developed and deployed S2SACC Outlook add-in using Yeoman and MS GraphAPI, facilitating document transfer to ContentCloud via email, with role-based access based on distribution groups, serving \textbf{1000+ users}.
                
                \item Developed S2EDS Outlook add-in used to send documents to enterprise servers from mailbox using ReactJs and Yeoman, serving \textbf{500+ users}.
            }
        \resumeItemListEnd
    \resumeSubHeadingListEnd

%--------PROGRAMMING SKILLS------------
\section{Skills}
    \resumeSubHeadingListStart
        \resumeSubItem
            {Languages} {Python, C/C++, JavaScript/TypeScript, Rust, Java, Assembly, SQL}
        \resumeSubItem
            {AI \& Agent Frameworks} {LangGraph, LangChain, LangSmith, OpenAI SDK, Anthropic Claude SDK, Vercel AI SDK, DeepEval, AWS Bedrock (Titan, Claude, multimodal embeddings), TensorFlow}
        \resumeSubItem
            {Backend \& Frontend} {FastAPI, SQLAlchemy, Alembic, Next.js, React, PostgreSQL, MERN, Electron, Yeoman}
        \resumeSubItem
            {Cloud \& DevOps} {AWS (EKS, OpenSearch, Lambda, S3, API Gateway, Secrets Manager), Kubernetes, Helm, Docker, Terraform, Terragrunt, Okta OIDC, Git, Mainframe (z/OS)}
 \resumeSubHeadingListEnd

%--------Achievements------------
 \section{Accolades}
  \resumeSubHeadingListStart
      \resumeSubItem{KVPY}
      {Recognized as a Kishore Vaigyanik Protsahan Yojana Scholar in 2017-2018}
    \resumeSubItem{Deloitte}
      {Honored with 3 Excellence Awards for development of the S2SACC Add-in and the OpenText Server upgrade}
    \resumeSubItem{Warner Bros. Discovery}
      {Honored with an Excellence Award for designing and building Knowledge Hub}
  \resumeSubHeadingListEnd
%-------------------------------------------
\end{document}
